% ===============================================================
%  Title Box Styling (tcolorbox + abstract/keywords storage)
% ===============================================================

% \usepackage[most]{tcolorbox}

% \tcbuselibrary{skins}
% \tcbuselibrary{breakable}
\usepackage{environ}  % 用来存 abstract 和 keywords

% ------------------------------
% 1. 存储 Abstract 内容
% ------------------------------
\NewEnviron{myabstract}{%
  \global\def\savedabstract{\BODY}%
}
\newcommand{\savedabstract}{Dexterous grasp generation aims to produce grasp poses that align with task requirements and human‑interpretable grasp semantics. However, achieving semantically controllable dexterous grasp synthesis remains highly challenging due to the lack of unified modeling of multiple semantic dimensions, including grasp taxonomy, contact semantics, and functional affordance. To address these limitations, we present OmniDexVLG, a multimodal, semantics‑aware grasp generation framework capable of producing structurally diverse and semantically coherent dexterous grasps under joint language and visual guidance.
Our approach begins with OmniDexDataGen, a semantic‑rich dexterous grasp dataset generation pipeline that integrates grasp‑taxonomy–guided configuration sampling, functional‑affordance contact point sampling, taxonomy‑aware differential force‑closure grasp sampling, and physics‑based optimization and validation, enabling systematic coverage of diverse grasp types. We further introduce OmniDexReasoner, a multimodal grasp‑type semantic reasoning module that leverages multi‑agent collaboration, retrieval‑augmented generation (RAG), and chain‑of‑thought (CoT) reasoning to infer grasp‑related semantics and generate high‑quality annotations that align language instructions with task‑specific grasp intent. Building upon these components, we develop a unified Vision-Language-Grasping (VLG) generation model that explicitly incorporates grasp taxonomy, contact structure, and functional affordance semantics, enabling fine‑grained control over grasp synthesis from natural language instructions.
Extensive experiments in simulation and real‑world object grasping and ablation studies demonstrate that our method substantially outperforms state‑of‑the‑art approaches in terms of grasp diversity, contact semantic diversity, functional affordance diversity, and semantic consistency. These results highlight the critical role of multi‑dimensional semantic modeling, including grasp taxonomy, contact semantics, and affordance reasoning, in advancing dexterous grasp generation. More details are available on our project website
% ~\url{https://sites.google.com/view/omnidexvlg-anonymous}.
~\url{https://sites.google.com/view/omnidexvlg}.
}  % 默认空

% ------------------------------
% 2. 存储 Keywords 内容
% ------------------------------
\NewEnviron{mykeywords}{%
  \global\def\savedkeywords{\BODY}%
}
\newcommand{\savedkeywords}{
LLM/VLM, Multi-Fingered Robotic Hand, Grasp Taxonomy, Generation Model.
}  % 默认空

% ===============================================================
% 3. 彩色标题框
% ===============================================================

\makeatletter

\newcommand{\mycolortitlebox}{%
  \begin{tcolorbox}[
    enhanced,
    breakable,
    % colback=blue!3,           % 背景颜色
    % colframe=blue!60!black,   % 边框颜色
    colback=green!3,           % 背景颜色
    colframe=green!60!black,   % 边框颜色
    boxrule=0.9pt,
    % sharp corners=downhill,
    arc=6pt,
    arc=10pt,
    left=6mm,right=6mm,top=6mm,bottom=6mm
  ]

    % ---------- Title & Author ----------
    \begin{center}
      {\LARGE\bfseries \@title \par}
      \vskip 0.8em
      {\large \@author \par}
    \end{center}

    \vskip 1em

    % ---------- Abstract ----------
    {\bfseries Abstract — } \savedabstract\par

    \vspace{1em}

    % ---------- Index Terms ----------
    {\bfseries Keywords — } \savedkeywords\par

  \end{tcolorbox}
}

% ===============================================================
% 4. 重写 \maketitle：双栏模式跨栏置顶
% ===============================================================

\renewcommand{\maketitle}{%
  \begingroup
    \normalfont

    \setcounter{footnote}{0}%
    \renewcommand{\thefootnote}{\fnsymbol{footnote}}%

    % -------- 双栏模式 --------
    \if@twocolumn
      \twocolumn[
        \begin{@twocolumnfalse}
          \mycolortitlebox
          \vspace{1.0em}
        \end{@twocolumnfalse}
      ]
    % -------- 单栏模式 --------
    \else
      \mycolortitlebox
    \fi

    % 恢复 footnote 格式
    \setcounter{footnote}{0}%
    \renewcommand{\thefootnote}{\arabic{footnote}}%

    \thispagestyle{empty}
  \endgroup
}

\makeatother
